\documentclass[12 pt, letterpaper]{hmcpset}
\usepackage{hw-preamble}

\name{Jayden Thadani}
\class{MATH 142 --- Differential Geometry}
\assignment{Homework assignment \#6}
\duedate{8th March 2024}

\begin{document}

\begin{problem}[Problem 1.]
	If $\mathbf{W}$ is a vector field of constant length $\lVert \mathbf{W} \rVert$, prove that for any vector field $\mathbf{V}$, the covariant derivative $\nabla_{\mathbf{V}} \mathbf{W}$ is everywhere orthogonal to $\mathbf{W}$.
\end{problem}

\begin{solution}
	Note that if we have $\mathbf{W} = [(w_{1}, w_{2}, w_{3}), \mathbf{p}]$, then we have $\sum_{i = 1}^{3} w_{i}^{2}(\mathbf{p})$ is constant always. Taking the gradient on both sides gives us
	\[
		\sum_{i = 1}^{3} w_{i}(\mathbf{p}) \nabla w_{i}(\mathbf{p}) = 0.
	\]
	Dotting with $\mathbf{v}$ from $\mathbf{V} = [\mathbf{v}, \mathbf{p}]$, we get
	\[
		\mathbf{v} \cdot \sum_{i = 1}^{3} w_{i}(\mathbf{p}) \nabla w_{i}(\mathbf{p}) = 0
	\]
	but we can rewrite this
	\[
		\begin{split}
			\mathbf{v} \cdot \sum_{i = 1}^{3} w_{i}(\mathbf{p}) \nabla w_{i}(\mathbf{p}) & = \sum_{i = 1}^{3} w_{i}(\mathbf{p}) (\nabla w_{i}(\mathbf{p}) \cdot \mathbf{v}) \\
												     & = \left ( \sum_{i = 1}^{3} w_{i}(\mathbf{p}) \, \mathbf{U}_{i}(\mathbf{p}) \right ) \cdot \left ( \sum_{i = 1}^{3} (\mathbf{v} \cdot \nabla w_{i}(\mathbf{p})) \, \mathbf{U}_{i}(\mathbf{p}) \right ) \\
												     & = \mathbf{W}(\mathbf{p}) \cdot \nabla_{\mathbf{V}} \mathbf{W} (\mathbf{p})
		\end{split}
	\]
	Thus, we have
	\[
		\begin{split}
			\mathbf{W}(\mathbf{p}) \cdot \nabla_{\mathbf{V}} \mathbf{W}(\mathbf{p}) & =  \mathbf{v} \cdot \sum_{i = 1}^{3} w_{i}(\mathbf{p}) \nabla w_{i}(\mathbf{p}) \\
												& = 0
		\end{split}
	\]
	or, in words, the covariant derivative $\nabla_{\mathbf{V}} \mathbf{W}$ is everywhere orthogonal to $\mathbf{W}$, as desired.
\end{solution}

\pagebreak

\begin{problem}[Problem 2.]
	Let $\mathbf{X} : \mathbb{R}^{3} \to \mathrm{T}\mathbb{R}^{3}$ be the special vector field such that $\mathbf{X}(\mathbf{p}) = [\mathbf{p}, \mathbf{p}]$. Prove that $\nabla_{\mathbf{V}} \mathbf{X} = \mathbf{V}$ for every vector field $\mathbf{V}$.
\end{problem}

\begin{solution}
	Note that since the function $\mathbb{R}^{3} \to \mathbb{R}^{3}$ by $\mathbf{p} \mapsto \mathbf{p}$ is linear, its Jacobian is just itself --- the identity linear map. Thus, for $\mathbf{X}(\mathbf{p}) = [\mathbf{x}, \mathbf{p}]$ with $\mathbf{x}(\mathbf{p}) = \mathbf{p}$ we have that the Jacobian of $\mathbf{x}$ is
	\[
		D \mathbf{x} \Big |_{\mathbf{p}} = I
	\]
	at every point $\mathbf{p}$, where $I$ is the identity linear map on $\mathbb{R}^{3}$.
	Since an alternate defintion of the covariant derivative $\nabla_{\mathbf{V}} \mathbf{X}$ is
	\[
		\nabla_{\mathbf{V}} \mathbf{X}(\mathbf{p}) = [D \mathbf{x} \Big |_{\mathbf{p}} (\mathbf{v}), \mathbf{p}]
	\]
	for any vector field $\mathbf{V} = [\mathbf{v}, \mathbf{p}]$ we have
	\[
		\begin{split}
			\nabla_{\mathbf{V}} \mathbf{X}(\mathbf{p}) & = [D \mathbf{x} \Big |_{\mathbf{p}} (\mathbf{V}(\mathbf{p})), \mathbf{p}] \\
								   & = [I(\mathbf{v}), \mathbf{p}] \\
								   & = [\mathbf{v}, \mathbf{p}] \\
								   & = \mathbf{V}(\mathbf{p}).
		\end{split}
	\]
	That is, we have $\nabla_{\mathbf{V}} \mathbf{X} = \mathbf{V}$ as desired.
\end{solution}

\pagebreak

\begin{problem}[Problem 3.]
	If $\mathbf{E}_{1}, \mathbf{E}_{2}, \mathbf{E}_{3}$ is a frame field and $\mathbf{W} = \sum_{i} f_{i} \, \mathbf{E}_{i}$ is a vector field, prove the \textit{covariant derivative formula}:
	\[
		\nabla_{\mathbf{V}} \mathbf{W} = \sum_{j = 1}^{3} \left ( \mathbf{V}[f_{j}] + \sum_{i = 1}^{3} f_{i} \, \omega_{ij}(\mathbf{V}) \right ) \mathbf{E}_{j}
	\]
	for any vector field $\mathbf{V}$.
\end{problem}

\begin{solution}
	Note that by the linearity and the Leibniz rule for the covariant derivative we have
	\[
		\begin{split}
			\nabla_{\mathbf{V}} \mathbf{W} & = \sum_{i = 1}^{3} \nabla_{\mathbf{V}} (f_{i} \, \mathbf{E}_{i}) \\
						       & = \sum_{i = 1}^{3} \mathbf{V}[f_{j}] \, \mathbf{E}_{i} + f_{i} \, \nabla_{\mathbf{V}}\mathbf{E}_{i} \\
						       & = \sum_{i = 1}^{3} \left ( \mathbf{V}[f_{i}] \, \mathbf{E}_{i} + f_{i} \sum_{j = 1}^{3} \omega_{ij}(\mathbf{V}) \, \mathbf{E}_{j} \right )
		\end{split}
	\]
	Notice that for some fixed $j \in \{ 1, 2, 3 \}$ each $\mathbf{E}_{j}$ is multiplied once by the coefficient $\mathbf{V}[f_{j}]$ and once by each of the coefficients $f_{i} \, w_{ij}$ for $i \in \{ 1, 2, 3 \}$. Thus, we can rearrange the sum into the desired form ---
	\[
		\nabla_{\mathbf{V}} \mathbf{W} = \sum_{j = 1}^{3} \left ( \mathbf{V}[f_{i}] + \sum_{i = 1}^{3} f_{i} \, \omega_{ij}(\mathbf{V}) \right ) \mathbf{E}_{j}.
	\]
\end{solution}

\pagebreak

\begin{problem}[Problem 4.]
	Let $\mathbf{E}_{1}, \mathbf{E}_{2}, \mathbf{E}_{3}$ be the cylindrical frame field. If $\mathbf{V}$ is a vector field such that $\mathbf{V}[r] = r$ and $\mathbf{V}[\vartheta] = 1$, compute $\nabla_{\mathbf{V}}(r \cos{\vartheta} \mathbf{E}_{1} + r \sin{\vartheta} \mathbf{E}_{2})$.
\end{problem}

\begin{solution}
	To compute this, again, we use linearity, the Leibniz rule, and the chain rule.
	\[
		\begin{split}
			\nabla_{\mathbf{V}}(r \cos{\vartheta} \, \mathbf{E}_{1} + r \sin{\vartheta} \, \mathbf{E}_{2}) & = \mathbf{V}[r \, \cos{\vartheta}] \, \mathbf{E}_{1} + \mathbf{V}[r \sin{\vartheta}] \, \mathbf{E}_{2} + r \cos{\vartheta} \, \nabla_{\mathbf{V}} \mathbf{E}_{1} + r \sin{\vartheta} \, \nabla_{\mathbf{V}} \mathbf{E}_{2} \\
														       & = (\mathbf{V}[r] \cos{\vartheta} + r \, \mathbf{V}[\cos{\vartheta}]) \mathbf{E}_{1} + (\mathbf{V}[r] \sin{\vartheta} + r \, \mathbf{V}[\sin{\vartheta}]) \mathbf{E}_{2} \\
														       & + r \cos{\vartheta} \, \omega_{12}(\mathbf{V}) \, \mathbf{E}_{2} - r \sin{\vartheta} \, \omega_{12}(\mathbf{V}) \, \mathbf{E}_{1} \\
														       & = (r \cos{\vartheta} - r \sin{\vartheta} \, \mathbf{V}[\vartheta]) \mathbf{E}_{1} + (r \sin{\vartheta} + r \cos{\vartheta} \, \mathbf{V}[\vartheta]) \mathbf{E}_{2} \\
														       & + r \cos{\vartheta} \, d\vartheta (\mathbf{V}) \, \mathbf{E}_{2} - r \sin{\vartheta} \, d\vartheta(\mathbf{V}) \, \mathbf{E}_{1} \\
														       & = (x - y (1 + d\vartheta(\mathbf{V}))) \mathbf{E}_{1}
		\end{split}
	\]
	Notice that $r \cos{\vartheta} = x$ and $r \sin{\vartheta} = y$. Then using linearity and the Leibniz rule we get
	\[
		\begin{split}
			\nabla_{\mathbf{V}}(r \cos{\vartheta} \, \mathbf{E}_{1} + r \sin{\vartheta} \, \mathbf{E}_{2}) & = \nabla_{\mathbf{V}}(x \, \mathbf{E}_{1} + y \, \mathbf{E}_{2}) \\
														       & = \mathbf{V}[x] \, \mathbf{E}_{1} + x \nabla_{\mathbf{V}} \mathbf{E}_{1} + \mathbf{V}[y] \, \mathbf{E}_{2} + y \nabla_{\mathbf{V}} \mathbf{E}_{2} \\
														       & = (v_{1} - y \, \omega_{12}(\mathbf{V})) \mathbf{E}_{1} + (v_{2} + x \, \omega_{12}(\mathbf{V})) \mathbf{E}_{2} \\
														       & = 
		\end{split}
	\]
\end{solution}

\pagebreak

\begin{problem}[Problem 5.]
	For a $1$-form $\phi = \sum_{i} f_{i} \, \theta_{i}$, prove
	\[
		d \phi = \sum_{j = 1}^{3} \left ( df_{j} + \sum_{i = 1}^{3} f_{i} \, \omega_{ij} \right ) \wedge \theta_{j}.
	\]
\end{problem}

\begin{solution}
	Note that by linearity and the Leibniz rule for exterior derivatives we can write
	\[
		\begin{split}
			d \phi & = \sum_{i = 1}^{3} d(f_{i} \, \theta_{i}) \\
			       & = \sum_{i = 1}^{3} df_{i} \wedge \theta_{i} + f_{i} \, d\theta_{i} \\
			       & = \sum_{i = 1}^{3} \left ( df_{i} \wedge \theta_{i} + f_{i} \sum_{j = 1}^{3} \omega_{ij} \wedge \theta_{j} \right ).
		\end{split}
	\]
	Notice that for some fixed $j \in \{ 1, 2, 3 \}$, each $\theta_{j}$ is wedged with $df_{j}$ once and once with each of the coefficients $f_{i} \, w_{ij}$ for $i \in \{ 1, 2, 3 \}$ (in all of the cases $\theta_{j}$ is on the right). Thus, by the linearity of the wedge product, we can rearrange the sum into the desired form ---
	\[
		d \phi = \sum_{j = 1}^{3} \left ( df_{i} + \sum_{i = 1}^{3} f_{i} \, \omega_{ij} \right ) \wedge \theta_{j}.
	\]
\end{solution}

\end{document}
